\section{Observations, assumptions, and forward quantities}
The observable is transverse displacement between direct S and core-reflected ScS.
The inversion assigns signed waveform amplitude to candidate scattering positions;
it does not recover an absolute shear-velocity perturbation or a temperature field.
We implement isotropic, single-scattering diffraction stacking in a spherical
reference Earth. Multiple scattering, elastic conversions, azimuthal anisotropy,
and finite-frequency triplication amplitudes are outside this imaging approximation.
The interpretation of a folded slab is kept separate from the calculated image.

The scientific specifications are the main text and supplementary legends of
Hutko et al. (2006), with processing details from Lay et al. (2004) and Avants
et al. (2006). The Nature legends take precedence where these descriptions differ.
They specify magnitudes 6.0--7.2, source depths 218--610 km, epicentral distances
65--85 degrees, and a final low-pass corner of 0.3 Hz. They normalize migration
traces to S; the earlier stacking study instead normalized to ScS.

We identify historical earthquake dates from Table 2 of Thomas et al. (2004).
That table contains 15 entries but repeats event number 7; the text says 14.
The entry dated 30 August 1999 is identified with the USGS event of 30 August
1998: its latitude and longitude match the printed values exactly and its depth
differs by only 0.5 km. We treat the printed year as a probable typographical
error and preserve this inference in the catalog. Its preferred magnitude is
5.7, a second historical exception to the stated selection bounds. The USGS
preferred magnitude of the identified 1991 event is 7.3, outside the published
range. Historical identities are preserved and such catalog exceptions are
flagged; the new-event query retains the written magnitude and depth bounds.
The historical 23 April 2000 entry is identified with the 17:01 UTC event,
consistent with the location in the published table, rather than its larger
09:27 UTC neighbour. No hypocentre used in computation is copied from that table.
Magnitudes are USGS preferred magnitudes of mixed types, retained in the event
tables. This is not guaranteed to duplicate the original CMT magnitude scale.

All event records and moment tensors are fetched from USGS. All waveforms and
StationXML responses are fetched from EarthScope, NCEDC, or SCEDC. Every request
records its URL, retrieval date, SHA-256 hash, size, and status; failures are saved
rather than silently filled. Natural Earth coastlines are used only for geographic
context. No author waveform collection, simulation output, code, or plotting
script is consulted. Published figure pages are used only for comparison, never
as numerical arrays or digitized model inputs.

Polet and Kanamori (2002) do report illustrative regional corrections: an
approximately east--west fast direction and 1.2 s delay at PAS, and a Bay-area
two-layer example with an upper 135-degree, 0.7 s layer above an east--west,
1.4 s layer. These are nonunique regional models; they do not identify the
complete station/path correction set used by Hutko et al. We do not silently
assign these examples to every station. Their existence means that the missing
correction set is a limitation of this reconstruction, not an assertion that
the literature contains no numerical anisotropy constraints.

The historical station list is an explicit approximation to California broadband
coverage, augmented by named western US stations. For later events the acquisition
also queries BK, CI, IU, II, US, and TA stations between 30--49 degrees N and
126--105 degrees W. Historical sites are retained; additional stations are sampled
at one site per network per one-degree cell to limit redundant downloads. This
thinning is a new acquisition assumption, not a published selection rule. The
expanded analysis includes the retained historical observations. The same-region
expanded subset retains every accepted historical trace and requires newly added
traces to have calculated ScS bounce points within 0--15 degrees N and
93--87 degrees W. The wider analysis uses all retained paths and masks plots with
fewer than 20 contributing traces or three events. These thresholds describe
illumination, not guaranteed spatial resolution.

\subsection{Instrument correction, polarization, and source removal}
Each gap-free BH three-component record is detrended, tapered over 3\% of its
length, and deconvolved to displacement using the epoch-appropriate archive
response. The response prefilter corners are 0.003, 0.005, 4, and 5 Hz; the upper
corners are reduced if required by the input Nyquist frequency. The instrument
water level is 60 dB. Sensor orientations from StationXML determine the rotation
to north and east. A fourth-order two-pass 0.01 Hz high-pass filter and an
eighth-order two-pass 2 Hz antialias filter precede sampling at 5 Hz. These fixed
choices replace the incompletely documented trace-dependent filters in the papers.

For backazimuth $b$, measured clockwise from north at the station, the transverse
component is
\begin{equation} u_T=u_N\sin b-u_E\cos b. \end{equation}
An optional receiver correction rotates to a fast direction $\phi$, delays the
fast component by $\delta t/2$, and advances the slow component by $\delta t/2$,
then rotates back. This preserves the average arrival time. The correction is
implemented and tested, but the supplied numerical correction file is empty:
the allowed written descriptions do not provide a complete usable station table,
and published splitting maps were not digitized. Therefore the observational
images in this report have \emph{no applied receiver splitting correction}.
This omission is a material difference from the publication.

USGS moment tensors are transformed from the spherical up--south--east basis to
north--east--down. The predicted SH radiation coefficient is
\begin{equation} R_{SH}=\frac{\boldsymbol{e}_T^{\mathsf T}\boldsymbol{M}\boldsymbol{n}}{M_0},
\qquad M_0=\left(\tfrac12\sum_{ij}M_{ij}^2\right)^{1/2}, \end{equation}
where $\boldsymbol{n}$ is the takeoff direction and $\boldsymbol{e}_T$ the transverse
unit vector. Where tensors are available, S and ScS must have matching signs and
$|R_{SH}|\geq0.6$. Unavailable tensors are flagged rather than invented; accepted
unverified records are reported separately. A coefficient range reported for the
original CMT solutions need not hold for the current USGS tensor, so this check
can remove an originally used event.

Phase peaks are measured on the 0.3 Hz low-pass displacement. S is sought within
15 s of its PREM prediction. Its residual defines a common time shift for the
ScS search; ScS is sought within 15 s of that shifted prediction and at least
6 s after S, with the S polarity. Picks on a window boundary are rejected.
This automated rule replaces unpublished manual picks and can confuse interfering
phases. We require a minimum peak-to-RMS amplitude ratio of 3 for both phases,
using the interval 90--30 s before picked S for noise. Colocated sensors appearing
under different network codes are represented once. All choices and rejections
remain in the trace audit.

For each event, up to ten candidates with the highest signal-to-noise ratios are
aligned on ScS and normalized by its signed amplitude. At least three candidates
are required. Their mean over $[-8,14]$ s is cropped at the nearest zero crossings
around its principal peak and tapered with a 25\% Tukey taper. This isolates an
empirical source wavelet $w$ and limits inclusion of precursor energy. It is an
assumed reproducible replacement for the original manual choice of simple ScS
pulses. A minimum correlation of 0.65 with this wavelet is required.

For displacement spectrum $U$ and wavelet spectrum $W$, the source-deconvolved
signal is
\begin{equation}
 D(\omega)=\frac{U(\omega)W^*(\omega)}{
 \max\{|W(\omega)|^2,\epsilon\max_{\omega'}|W(\omega')|^2\}}\,G(\omega),
 \qquad \epsilon=10^{-4}.
\end{equation}
Here $G$ is a fourth-order, two-pass Butterworth low-pass filter with corner
0.3 Hz. The FFT is zero padded beyond the combined signal and source lengths;
the wavelet zero time is explicitly placed at index zero. The deconvolved S and
ScS peaks are repicked within 5 s, and $d(t)$ is divided by its signed S peak.
The event supplements show both pre-deconvolution displacement and deconvolved
signals, each separately S-normalized, with at least 80 s before S and 65 s
after ScS. Display clipping at twice the S amplitude is labelled; saved traces
are not clipped. These diagnostics do not replace expert manual acceptance.

\subsection{Spherical ray integration and migration}
The Earth radius is $R=6371$ km and PREM's CMB depth is 2891 km. We use the
isotropic tabulation of PREM supplied with ObsPy. For ray parameter $p$ in
seconds per radian, shear speed $v(r)$, and a monotonically upgoing leg,
\begin{align}
 \Delta(p)&=\int_{r_x}^{r_e}
 \frac{p\,v(r)}{r^2\sqrt{1-[p v(r)/r]^2}}\,dr,\\
 T(p)&=\int_{r_x}^{r_e}
 \frac{1}{v(r)\sqrt{1-[p v(r)/r]^2}}\,dr.
\end{align}
The endpoint $e$ is a source or receiver and $x$ a candidate scatterer.
Only $p<\min r/v(r)$ is admitted; legs that would turn below $x$ are absent.
We split radial integration at every material boundary, endpoint, and 15 km
interval. Sixteen-point Gauss--Legendre quadrature uses $r-r_x=u^2$ to regularize
near-grazing incidence. One thousand ray parameters, clustered near the upper
limit, define monotone travel-time curves. Linear interpolation onto a
0.02-degree distance table avoids separate ray tracing for every voxel.
Direct S, ScS, and multiple predictions used for picks come from ObsPy TauP;
the scattering-leg integrator is independently implemented.

The original image domain is $[-6,20]$ degrees latitude, $[-103,-77]$ degrees
longitude, and 2310--2910 km depth, at 0.2-degree horizontal and 10 km vertical
sampling. For the last 19 km below the CMB, basal mantle shear speed is extended
solely to map post-ScS time. It is not interpreted as a shear wave in the fluid
core. The expanded domain is $[-10,25]$ degrees latitude and $[-110,-70]$ degrees
longitude at the same spacing. Geographic latitudes are treated as spherical;
ellipticity, topography, and station elevation travel-time corrections are omitted.
Profile endpoints are specified in configuration, independently of the published
plots. Horizontal distances are measured at the CMB radius, with no vertical
exaggeration. This differs from measuring profile length at the Earth's surface.

For reference phase $q$ (S or ScS), the sampled time is
\begin{equation}
 t_{iq}(\boldsymbol{x})=t^{\rm pick}_{iq}+
 T(s_i,\boldsymbol{x})+T(\boldsymbol{x},r_i)-T^{\rm PREM}_{iq}.
\end{equation}
The measured reference establishes a datum that reduces origin-time and shared
path errors. Direct S and predicted sS and SS are multiplied by symmetric
raised-cosine notches $m(t)=\sin^2\{(\pi/2)\min(|t-t_0|/2\,\mathrm{s},1)\}$.
The ScS-censored variants apply the same 4 s full-width notch to ScS. The exact
published taper shape is not specified. Our hitcount counts valid paths and
samples, including samples with tapered amplitudes, rather than a count of
nonzero amplitudes. Consequently its retained/censored versions can coincide.

For $N(\boldsymbol{x})$ valid traces, both the arithmetic mean
$\bar A=\sum_i d_i[t_{iq}(\boldsymbol{x})]/N(\boldsymbol{x})$ and the displayed
hitcount-weighted mean are saved:
\begin{equation}
 A(\boldsymbol{x})=\bar A(\boldsymbol{x})\,
 \frac{N(\boldsymbol{x})}{\max_{\boldsymbol{x}'}N(\boldsymbol{x}')}.
\end{equation}
This explicit interpretation of ``weighted according to hitcount'' is an
assumption. No geometrical-spreading or absolute scattering-amplitude inversion
is performed. The display saturates at $\pm0.15$ in S-normalized amplitude units;
unsaturated arrays remain available. No positive reflector or step is imposed
on observational images.

The first-order tomography correction operator is
\begin{equation} \delta T=-\int_{\rm ray}\frac{\delta v}{v}\,dt.\end{equation}
Its differential contribution is the sum on the two scattering legs minus
that on the reference ray. The segment operator is implemented and verified for
a uniform perturbation. The UT, UCB, and CIT coefficients are not specified by
the written descriptions. We do not manufacture these models, download the
authors' numerical inputs, or substitute generic velocity anomalies under their
names. These three Figure 3 panels remain unreproduced.

\subsection{Additional tests and the boundary of the reconstruction}
The double-array test solves $\Delta_s(p)+\Delta_r(p)=\Delta_{sr}$ at each
candidate reflector depth and samples each trace at that reflection time.
ScS bounce latitude 6 degrees N separates southern and northern bins. Confidence
intervals resample entire events 300 times, preserving within-event correlation.
They are pointwise intervals, not corrections for searching across depths.
The source-location sensitivity test perturbs only sources south of 25 degrees S
by $\pm40$ km in depth or by 40 km toward/away from the selected profile midpoint.

Synthetic point tests inject one-second Gaussian arrivals from two explicitly
assumed isotropic scatterers at 40\% and 70\% along profile A, at depths
2611 and 2711 km (100 km vertical separation). These test the
diffraction-stack operator and acquisition geometry, not full elastic propagation.
Each pulse has unit positive peak and a Gaussian standard deviation of 1 s;
attenuation, conversions, and additional reflections are absent. The same two
scatterers are tested with each dataset's own accepted event/station geometry.
A separate kinematic SH test uses a 2.6\% shear-speed increase 264 km above the
CMB and optional 2.6, 5, or 10\% decreases 90 km above it. With equal density,
its local plane-interface displacement coefficient is
\begin{equation}
 \mathcal R_{SH}=\frac{\mu_1 q_1-\mu_2 q_2}{\mu_1 q_1+\mu_2 q_2},\qquad
 q_j=\sqrt{v_j^{-2}-p_h^2},\quad \mu_j=\rho v_j^2.
\end{equation}
Complex $q_j$ permits a postcritical phase shift. This illustrative ray calculation
does not include the refracted Scd branch, caustic amplitudes, attenuation, or
full-wave reflectivity. It therefore cannot reproduce the published synthetic
waveform panels. The 1.5\% step model at 170/270 km above the CMB is generated
as an input-model illustration on PREM. No 2.5-D finite-difference solver or
thermomechanical slab-flow calculation is claimed.

There are no time-evolution initial conditions in the observational migration:
its input is the measured, processed time series, zero outside its sampled
support. The ray reference has a fluid-core SH reflection boundary and a solid
surface as encoded in PREM/TauP. The kinematic examples assume causal localized
impulses and isolated interface interactions; no absorbing numerical boundary or
dynamic initial stress field is applicable to them. A full-wave reconstruction
would require these additional specifications and validation.
